\section{Introduction}
\label{sec:introduction}

Retrieval-augmented generation (RAG) is the dominant deployment
pattern for large language models that must ground their outputs in
external evidence \citep{lewis2020rag,izacard2020fid,gao2023rag}.
The core empirical question --- \emph{when is a generated answer
faithful to its retrieved context?} --- has drawn an expanding
catalogue of mitigations: corrective retrieval
\citep{yan2024crag}, self-reflective generation
\citep{asai2024selfrag}, hierarchical summarization
\citep{sarthi2024raptor}, and selective per-passage refinement
\citep{xu2023recomp,jiang2023active}. The relative merit of these
methods is typically claimed via faithfulness measurements on a
small set of generators and one entailment scorer, often with
hyperparameters tuned on a held-out slice of the same dataset that
is then used to report results.

This paper examines the consequences of those evaluation choices.
We adopt one recently-named phenomenon ---
the \emph{refinement paradox}, in which more aggressive
per-passage refinement is reported to \emph{reduce} faithfulness ---
as a stress test. We pre-register six experiments across five
weaknesses and run them on free compute. We report what holds and
what does not, with effect sizes and confidence intervals throughout.

\paragraph{What we find.}
The headline magnitude of the refinement paradox does not survive
proper-power evaluation. At \(n{=}500\) and five seeds, only 1/5
seeds reaches the conventional \(p<0.05\) threshold, and the pooled
bootstrap 95\,\% confidence interval on the paired contrast crosses
zero. A pre-registered matched-similarity causal intervention finds
no support for context coherence as a causal mediator at fixed
mean retrieval similarity. Faithfulness scoring is itself unstable
across scorers: under DeBERTa-v3 the paradox magnitude is
0.011; under a RAGAS-style LLM-judge using a separate Mistral-7B
instance it is 0.140 ($\sim 13\times$ larger), with pairwise Pearson
\(r\) of 0.18 between the two on \(n{=}7{,}500\) paired evaluations.
Hyperparameter generalization is uneven: SQuAD-tuned thresholds do
not transfer to PubMedQA or NaturalQuestions without measurable
drop. Head-to-head comparison of four retrieval-time interventions
shows no clear winner --- CRAG wins HotpotQA, RAPTOR-2L narrowly
edges HCPC-v2 on SQuAD faithfulness, and Self-RAG underperforms
substantially when the evaluation harness is held fixed.

\paragraph{What survives.}
A structural disentanglement experiment shows a non-trivial
positive finding: at matched per-passage similarity,
\emph{coherence-preserving uninformative noise} (passages from the
same topic that do not contain the answer) produces a smaller
faithfulness drop than random off-topic noise. The two slopes are
\(-0.043\) and \(-0.069\) per noise rate, respectively. Context
coherence appears to carry some signal independent of similarity ---
even when the canonical operationalization (CCS, the Context
Coherence Score) is not the causal mechanism the matched-similarity
intervention tests.

\paragraph{Contributions.}
This paper makes four contributions.
\begin{enumerate}
\itemsep -2pt
\item A pre-registered audit of the refinement paradox at
$10\times$ the published sample size (5 seeds, 7{,}500 paired
contrasts), showing the published magnitude collapses at scale.
\item A multi-metric triangulation of RAG faithfulness scoring on
$n{=}7{,}500$ paired generations, revealing that DeBERTa-v3 NLI is
an outlier relative to roberta-large-mnli and an LLM-judge.
\item A pre-registered causal intervention with matched-similarity
HIGH/LOW CCS triples that does not support a causal interpretation
of context coherence, alongside a coherence-preserving noise
experiment that nonetheless preserves a structural role for
coherence beyond similarity.
\item A released benchmark
(\href{https://huggingface.co/datasets/saketmgnt/context-coherence-bench}{\texttt{saketmgnt/context-coherence-bench}};
DOI \href{https://doi.org/10.5281/zenodo.19757291}{10.5281/zenodo.19757291}),
with per-query CSVs, pre-registrations, and Kaggle T4\(\times\)2
notebooks that reproduce every claim in this paper.
\end{enumerate}

\paragraph{Reading guide.}
\S\ref{sec:background}--\S\ref{sec:methodology} fix definitions and
the pre-registration discipline.
\S\ref{sec:causal_intervention}--\S\ref{sec:headtohead} report the
five experiments described above, each as a self-contained
subsection with hypothesis, sample, statistical test, and result.
\S\ref{sec:benchmark} documents the released artifact;
\S\ref{sec:related}--\S\ref{sec:conclusion} contextualize the work
and discuss limitations.
